\subsection{Скорость сходимости методов с предобуславливанием и затуханием весов}

Хотя шаг оптимизации весов модели может показаться простым, на него можно посмотреть с другой стороны. Вынесем матрицу $D_t^{-1}$ за скобки, что даёт нам следующий шаг:

\begin{equation}
\label{eq:wd_step}
w_{t+1} = w_t - \eta \cdot D_t^{-1} (\nabla f(w_t) + D_t \nabla r(w_t)).
\end{equation}

Это подталкивает нас к тому, чтобы ввести новую функцию $\widetilde{r}_t$, такую, что $\nabla \widetilde{r}_t(w) = D_t \nabla r(w)$, и новую целевую функцию
\begin{equation} \label{F_tilde}
    \widetilde{F}_t(w) := f(w) + \widetilde{r}_t(w),
\end{equation}
где новая целевая функция $\widetilde{F}_t$ меняется на каждом оптимизационном шаге, так как $D_t$ тоже обновляется на каждой итерации.

Новый адаптивный регуляризатор $\widetilde{r}_t$ в общем случае, к сожалению, не существует. Поэтому мы наложим ограничения на исходный регуляризатор и структуру предобуславливателя, которые оформлены в виде следующих предположений.

\begin{assumption}{(Структура регуляризатора)}
    \label{ass:regstruct}
    Регуляризатор $r$ сепарабелен, то есть он может быть представлен в следующем виде:
    $$r(w) = \sum_{i=1}^d r_i(w^i),$$
    где $r_i(x) \ge 0$ для $i \in \overline{1, d}$ и $x \in \mathbb{R}$.
\end{assumption}

\begin{assumption}{(Структура матрицы предобуславливания)}
    \label{ass:precondstruct}
    Матрица предобуславливания $D_t$ диагональна:
    $$ D_t = \textrm{diag} \left\{ d_t^1, \ldots, d_t^d \right\}.$$
\end{assumption}

Хотя эти предположения являются достаточно сильными, они выполняются для упомянутых ранее методов с предобуславливанием и затуханием весов, а также для таких популярных функций регуляризации, как регуляризация Тихонова и (сглаженная) LASSO-регуляризация.
Скорость сходимости обычно измеряется количеством итераций, которое необходимо для достижения заданного уровня погрешности. Чтобы получить такие оценки, мы должны ввести определённые предположения об оптимизируемой функции потерь. На протяжении всего последующего анализа мы предполагаем, что функции $f : \mathbb{R}^d \rightarrow \mathbb{R}$ и $r : \mathbb{R}^d \rightarrow \mathbb{R}$ являются гладкими в следующем смысле.

\begin{assumption}{($L$-гладкость)}
\label{ass:smoothness}
\begin{itemize}
    \item 	Функция $f$ является $L_f$-гладкой, то есть её градиент липшицев: существует константа $L_f > 0$ такая, что $\forall x, y \in \mathbb{R}^d$
    	\begin{equation*}
    		\|\nabla f(x) - \nabla f(y)\| \leq L_f \|x - y\|.
    	\end{equation*}
    \item    Функция $r$ является $L_r$-гладкой: существует константа $L_r > 0$ такая, что $\forall x, y \in \mathbb{R}^d$
	\begin{equation*}
		\|\nabla r(x) - \nabla r(y)\| \leq L_r \|x - y\|.
	\end{equation*}
\end{itemize}
\end{assumption}

Из липшицевости градиента стандартным образом следует квадратичная верхняя оценка (лемма о спуске), которой мы будем пользоваться в доказательствах: для $L$-гладкой функции $h$ и любых $x, y \in \mathbb{R}^d$
\begin{equation}
\label{eq:descent_lemma}
    h(x) \leq h(y) + \langle \nabla h(y), x-y \rangle + \frac{L}{2} \|x - y\|^2.
\end{equation}
Отметим, что то, что $f$ дважды дифференцируема, в доказательствах не используется; это нужно лишь для корректности предобуславливателей гессианного типа (OASIS, AdaHessian), в которых новая информация для матрицы предобуславливания строится по второй производной.

Для работы в невыпуклом случае нам понадобится ограничение на значения функции регуляризации; оно сформулировано в предположении \ref{ass:regbound}.

\begin{assumption}{(Ограниченность регуляризатора)}
\label{ass:regbound}
Регуляризатор ограничен: существует константа $\Omega > 0$ такая, что $\forall w \in \mathbb{R}^d$
\begin{equation*}
|r(w)| \le \Omega.
\end{equation*}
\end{assumption}

\begin{note}
Предположение \ref{ass:regbound} исключает неограниченные на всём пространстве регуляризаторы (например, регуляризацию Тихонова $r(w) = \frac{\lambda}{2}\|w\|^2$). Оно используется только в невыпуклом случае (теоремы \ref{theor:1} и \ref{theor:3}) для контроля дрейфа $\widetilde{F}_t$ по $t$, причём достаточно, чтобы оно выполнялось на множестве, содержащем траекторию метода: если траектория $\{w_t\}$ остаётся в шаре радиуса $R_{\max}$, можно положить $\Omega = \max_{\|w\| \le R_{\max}} r(w)$. В сильно выпуклом случае (теоремы \ref{theor:2} и \ref{theor:4}) это предположение не требуется, и анализ покрывает регуляризацию Тихонова без оговорок.
\end{note}

Мы используем стандартное ограничение на матрицу предобуславливания, которое сформулировано в предположении \ref{ass:preconditioned}.

\begin{assumption}{(Ограниченность предобуславливателя)}
\label{ass:preconditioned}
Существуют константы $\alpha, \Gamma \in \mathbb{R} : 0 < \alpha < \Gamma$ такие, что
\begin{equation*}
\alpha I \preccurlyeq D_t \preccurlyeq \Gamma I \Leftrightarrow \frac{I}{\Gamma} \preccurlyeq D_t^{-1} \preccurlyeq \frac{I}{\alpha}.
\end{equation*}
\end{assumption}

В работе~\citep{beznosikov2022scaled} (в контексте седловых задач), а также в работах~\citep{jahani2021doubly, sadiev2022stochastic} показано, что это предположение выполняется для популярных алгоритмов с предобуславливанием, таких как Adam, RMSProp и OASIS (для Adagrad --- при дополнительном условии ограниченности градиентов).


В нашем анализе мы рассматриваем два способа обновления матрицы предобуславливания.
В первом методе матрица обновляется через квадраты:
\begin{equation}
\label{eq:squares}
    (D_{t+1})^2 = \beta (D_t)^2 + (1 - \beta) (H_t)^2,
\end{equation}
где $H_t$ --- диагональная матрица, содержащая новую информацию, а $\beta \in [0, 1]$ --- параметр импульса (в практических методах $\beta \in (0, 1)$ и близок к 1).
Этот подход используется в Adam, а также в более ранних методах, таких как RMSProp и AdaHessian.
Второй способ предполагает использование первых степеней матриц:
\begin{equation}
\label{eq:linear}
    D_{t+1} = \beta D_t + (1 - \beta) H_t.
\end{equation}

Этот подход используется в OASIS --- недавно предложенном методе. В обоих случаях параметр импульса $\beta$ обычно выбирается близким к 1, что означает, что $D_t$ незначительно меняется в ходе обучения; формально это утверждение сформулировано в лемме \ref{lemma:Dt}.

Выполнение предположения \ref{ass:preconditioned} имеет решающее значение для сходимости и теоретического анализа, поэтому в алгоритмах используется принудительное ограничение снизу диагонали матрицы $D_t$:
\begin{equation}
\label{eq:alpha}
    \hat{D}_{t+1}^{ii} = \max \{ \alpha, |D_{t+1}^{ii}| \}.
\end{equation}
Срезка по модулю согласована с алгоритмом OASIS, в котором оценка диагонали гессиана может иметь отрицательные элементы. Далее всюду, где не оговорено противное, под $D_t$ понимается именно срезанная матрица $\hat{D}_t$; таким образом, предположение \ref{ass:preconditioned} выполняется автоматически, если $|H_t^{ii}| \le \Gamma$ и $\alpha I \preccurlyeq \hat{D}_0 \preccurlyeq \Gamma I$.


\begin{lemma}
\label{lemma:Dt}
{(Эволюция $D_t$)}
    Предположим, что для начальной матрицы $D_0$ выполнены предположения \ref{ass:precondstruct} и \ref{ass:preconditioned}, матрицы $H_t$ диагональны и $|H_t^{ii}| \le \Gamma$ для всех $i$ и $t$, а $D_t$ эволюционирует в соответствии с \eqref{eq:squares}, \eqref{eq:alpha} или
    \eqref{eq:linear}, \eqref{eq:alpha}. Тогда справедливы следующие утверждения:

    \begin{enumerate}
        \item предположения \ref{ass:precondstruct} и \ref{ass:preconditioned} выполняются для $\hat{D}_t$ при всех $t$;
        \item $\|\hat{D}_{t+1} - \hat{D}_t\|_\infty \le \frac{(1 - \beta) \Gamma^2}{2\alpha} =: \rho_\beta$ для \eqref{eq:squares};
        \item $\|\hat{D}_{t+1} - \hat{D}_t\|_\infty \le 2(1 - \beta) \Gamma =: \rho_\beta$ для \eqref{eq:linear}.
    \end{enumerate}
\end{lemma}

Эта лемма доказана в Приложении \ref{appendix:lemmas}; при доказательстве мы опираемся на работу~\citep{beznosikov2022scaled}. Всюду далее $\rho_\beta$ обозначает соответствующую правую часть: $\rho_\beta = \frac{(1-\beta)\Gamma^2}{2\alpha}$ для \eqref{eq:squares} и $\rho_\beta = 2(1-\beta)\Gamma$ для \eqref{eq:linear}. Подчеркнём, что $\rho_\beta \to 0$ при $\beta \to 1$.

С помощью предположений \ref{ass:regstruct} и \ref{ass:precondstruct} мы можем доказать существование $\widetilde{r}_t$ и, следовательно, $\widetilde{F}_t$. Мы оформим это в лемме \ref{lemma:existence}. Функция $\widetilde{r}_t$ определена однозначно с точностью до аддитивной константы, и для наших оценок этого достаточно.

\begin{lemma}
\label{lemma:existence}
{(Существование $\widetilde{r}_t$)}
    Пусть выполняются предположения \ref{ass:regstruct} и \ref{ass:precondstruct}. Тогда функция $\widetilde{r}_t$ существует и имеет следующий вид:
    $$\widetilde{r}_t(w) = \sum_{i=1}^d d_t^i \, r_i(w^i).$$
\end{lemma}

Используя предположение \ref{ass:smoothness}, мы можем гарантировать гладкость $\widetilde{r}_t$ и оценить константу Липшица её градиента; это сформулировано в лемме \ref{lemma:tildesmoothness}.

\begin{lemma}\label{lemma:tildesmoothness}{($L$-гладкость $\widetilde{r}_t$)}
Пусть выполняются предположения \ref{ass:regstruct}, \ref{ass:precondstruct}, \ref{ass:smoothness}, \ref{ass:preconditioned}. Тогда функция $\widetilde{r}_t$ является $L_{\tilde{r}}$-гладкой с $L_{\tilde{r}} = \Gamma L_r$:
\begin{equation*}
    		\|\nabla \widetilde{r}_t(x) - \nabla \widetilde{r}_t(y)\| \leq L_{\tilde{r}} \|x - y\| \quad \forall x, y \in \mathbb{R}^d.
\end{equation*}
В частности, для $\widetilde{r}_t$ выполняется неравенство \eqref{eq:descent_lemma} с константой $L_{\tilde r}$, а функция $\widetilde{F}_t$ является $\widetilde{L}$-гладкой с $\widetilde{L} = L_f + \Gamma L_r$.
\end{lemma}

Чтобы получить дополнительные оценки сходимости, в сильно выпуклом случае мы накладываем предположение \ref{ass:muconvex} на целевую функцию потерь.
\begin{assumption}{(Сильная выпуклость)}
    \label{ass:muconvex}
    Существует константа $\mu_f > 0$ такая, что $\forall x, y \in \mathbb{R}^d$ выполняется:
    $$
    f(y) \geq f(x) + \langle \nabla f(x), y-x \rangle + \frac{\mu_f}{2} \|x-y\|^2.
    $$
\end{assumption}

\subsubsection{Детерминированный случай}

Сначала рассмотрим детерминированную (полноградиентную) постановку: $g_t = \nabla f(w_t)$. Используя введённые предположения, мы доказываем сходимость методов с предобуславливанием и затуханием весов. Наши результаты сформулированы в теоремах \ref{theor:1} и \ref{theor:2}; доказательства можно найти в Приложении \ref{appendix:theorems}.

\begin{theorem}[невыпуклый случай]
    \label{theor:1}
    Пусть выполняются предположения \ref{ass:regstruct}, \ref{ass:precondstruct}, \ref{ass:smoothness}, \ref{ass:regbound}, \ref{ass:preconditioned}, метод \eqref{eq:wd_step} использует точный градиент ($g_t = \nabla f(w_t)$), а матрица предобуславливания обновляется в соответствии с условиями леммы \ref{lemma:Dt}. Пусть шаг обучения удовлетворяет условию
    \begin{equation*}
        \eta < \frac{2 \alpha}{\widetilde{L}}, \qquad \widetilde{L} = L_f + \Gamma L_{r},
    \end{equation*}
    где $L_f, L_{r}$ --- константы гладкости функций $f$ и $r$ (константы Липшица их градиентов).
    Обозначим $\Delta_0 = \widetilde{F}_0(w_0) - f^*$, где $w_0 \in \mathbb{R}^d$ --- начальная точка, $\widetilde{F}_t$ определена в \eqref{F_tilde}, а $f^*$ --- оптимальное значение задачи \eqref{eq:general}, и положим
    $$\delta := \Omega \rho_\beta = \begin{cases}
        \frac{(1 - \beta)\Gamma^2 \Omega}{2\alpha} & \text{для } \eqref{eq:squares}, \\
        2(1 - \beta)\Gamma \Omega  &  \text{для } \eqref{eq:linear}.
    \end{cases}$$
    Тогда для любой точности $\varepsilon$, удовлетворяющей условию
    \begin{equation*}
        \varepsilon > \frac{\delta\Gamma}{\eta - \frac{\widetilde{L}\eta^2}{2\alpha}},
    \end{equation*}
    число итераций, необходимое для достижения $\min\limits_{0 \le t \le T} \|\nabla\widetilde{F}_t(w_t)\|^2 \leq \varepsilon$, ограничено величиной
    \begin{equation*}
      T = \mathcal{O}\left( \frac{\Delta_0 \Gamma}{(\eta - \frac{\widetilde{L}\eta^2}{2\alpha}) \left( \varepsilon -\frac{\delta\Gamma}{\eta - \frac{\widetilde{L}\eta^2}{2\alpha}}\right)} \right).
\end{equation*}
\end{theorem}

\begin{note}
Метод гарантированно уменьшает квадрат нормы градиента изменённой функции $\widetilde{F}_t$ лишь до порога $\delta\Gamma/(\eta - \widetilde{L}\eta^2/(2\alpha))$, поэтому условие на $\varepsilon$ существенно. Поскольку $\delta \propto (1-\beta)$, порог может быть сделан меньше любой наперёд заданной точности, если выбрать $\beta$ достаточно близким к 1.
\end{note}

В сильно выпуклом случае мы доказываем линейную скорость сходимости. Подчеркнём принципиальный момент, согласующийся с леммой \ref{lemma:lowerbondF} ниже: метод с затуханием весов сходится не к решению $w^*$ исходной задачи \eqref{F_big}, а к решению $\widetilde{w}^*_t$ изменённой задачи --- минимизации $\widetilde{F}_t$. Именно расстояние до этой (движущейся с ростом $t$) точки и убывает линейно.

\begin{theorem}[сильно выпуклый случай]
\label{theor:2}
    Пусть выполняются предположения \ref{ass:precondstruct}, \ref{ass:smoothness}, \ref{ass:preconditioned}, \ref{ass:muconvex}, регуляризатор --- тихоновский: $r(w) = \frac{\lambda}{2}\|w\|^2$, $\lambda > 0$, метод \eqref{eq:wd_step} использует точный градиент ($g_t = \nabla f(w_t)$), а матрица предобуславливания обновляется в соответствии с условиями леммы \ref{lemma:Dt}. Обозначим:
    $$\widetilde{\mu} := \mu_f + \lambda\alpha, \qquad \widetilde{L} := L_f + \lambda\Gamma, \qquad \widetilde{w}^*_t := \argmin_{w \in \mathbb{R}^d} \widetilde{F}_t(w), \qquad \Omega_0 := \frac{\|\nabla f(0)\|}{\widetilde{\mu}},$$
    $$R_t^2 := \|w_t - \widetilde{w}^*_t\|^2_{D_t},$$
    при этом $\widetilde{w}^*_t$ существует, единственно и $\|\widetilde{w}^*_t\| \le \Omega_0$ для всех $t$.
    Пусть шаг обучения и параметр импульса удовлетворяют условиям
    \begin{equation*}
        0 < \eta \le \frac{\widetilde{\mu}\alpha^2}{\Gamma \widetilde{L}^2}, \qquad \rho_\beta \le \frac{\eta\widetilde{\mu}\alpha}{4\Gamma}.
    \end{equation*}
    Тогда для всех $T \ge 1$
    \begin{equation*}
        R_T^2 \le \left(1 - \frac{\eta\widetilde{\mu}}{4\Gamma}\right)^T R_0^2 + \frac{64\,\Gamma^3\lambda^2\Omega_0^2}{\eta^2\widetilde{\mu}^4}\,\rho_\beta^2.
    \end{equation*}
    В частности, для любого $\varepsilon > 0$, если $\beta$ выбран так, что второе слагаемое не превосходит $\varepsilon/2$, то $R_T^2 \le \varepsilon$ после
    \begin{equation*}
        T = \frac{4\Gamma}{\eta\widetilde{\mu}}\log\frac{2R_0^2}{\varepsilon} = \mathcal{O}\left(\frac{\Gamma}{\eta\widetilde{\mu}}\log\frac{R_0^2}{\varepsilon}\right) \text{ итераций.}
    \end{equation*}
\end{theorem}

\begin{note}
Второе слагаемое --- цена нестационарности предобуславливателя: цель $\widetilde{w}^*_t$ дрейфует вместе с $D_t$, и скорость этого дрейфа контролируется величиной $\rho_\beta \propto (1-\beta)$ из леммы \ref{lemma:Dt}. При $\beta \to 1$ (или при замороженной матрице $D_t \equiv D$) остаточный член исчезает и метод сходится к $\widetilde{w}^*$ линейно, вплоть до машинной точности.
\end{note}

\subsubsection{Стохастический случай}

Теперь откажемся от точного градиента и проведём анализ в стохастической постановке. Требования к стохастическому градиенту формализованы в следующем предположении.

\begin{assumption}{(Стохастический градиент)}
\label{ass:expectations}
Обозначим через $\mathcal{F}_t$ $\sigma$-алгебру, порождённую всей случайностью, реализованной строго до розыгрыша $g_t$ (в частности, $g_0, \ldots, g_{t-1}$ и случайностью, использованной при построении предобуславливателя, например векторами Радемахера $z_0, \ldots, z_t$). Матрица $D_t$ предсказуема, то есть $\mathcal{F}_t$-измерима (строится по информации, полученной строго до вычисления $g_t$); при этом $w_t$ также $\mathcal{F}_t$-измерима. Стохастические градиенты $g_t$ условно несмещены и имеют равномерно ограниченную условную дисперсию:
\begin{equation*}
\mathbb{E}\left[ g_t \mid \mathcal{F}_t \right] = \nabla f (w_t), \qquad \mathbb{E}\left[ \|g_t - \nabla f(w_t)\|^2 \mid \mathcal{F}_t \right] \leq \sigma^2.
\end{equation*}
\end{assumption}

\begin{note}
Предсказуемость $D_t$ соответствует <<запаздывающему>> варианту предобуславливателя и является стандартным приёмом в анализе адаптивных методов; для OASIS она выполняется автоматически, если вектор Радемахера $z_t$ разыгрывается до $g_t$ и независим от него.
\end{note}

Шаг метода принимает вид $w_{t+1} = w_t - \eta D_t^{-1} \widehat{g}_t$, где $\widehat{g}_t := g_t + D_t \nabla r(w_t)$; при этом в силу предположения \ref{ass:expectations} и $\mathcal{F}_t$-измеримости $w_t$ и $D_t$
$$\mathbb{E}[\widehat{g}_t \mid \mathcal{F}_t] = \nabla \widetilde{F}_t(w_t), \qquad \mathbb{E}\left[\|\widehat{g}_t - \nabla \widetilde{F}_t(w_t)\|^2 \mid \mathcal{F}_t\right] \le \sigma^2,$$
то есть $\widehat{g}_t$ --- несмещённая стохастическая оценка градиента изменённой функции с той же дисперсией.

\begin{theorem}[стохастический невыпуклый случай]
    \label{theor:3}
    Пусть выполняются предположения \ref{ass:regstruct}, \ref{ass:precondstruct}, \ref{ass:smoothness}, \ref{ass:regbound}, \ref{ass:preconditioned}, \ref{ass:expectations}, а матрица предобуславливания обновляется в соответствии с условиями леммы \ref{lemma:Dt}. Пусть шаг обучения удовлетворяет условию
    \begin{equation*}
        0 < \eta \le \frac{\alpha^2}{\Gamma \widetilde{L}}, \qquad \widetilde{L} = L_f + \Gamma L_r.
    \end{equation*}
    Тогда в обозначениях теоремы \ref{theor:1}
    \begin{equation*}
        \min_{0 \le t \le T} \mathbb{E}\|\nabla \widetilde{F}_t (w_t) \|^2 \;\le\; \frac{2\Gamma\Delta_0}{\eta (T+1)} \;+\; \frac{\Gamma \widetilde{L} \eta \sigma^2}{\alpha^2} \;+\; \frac{2\Gamma\delta}{\eta}.
    \end{equation*}
\end{theorem}

\begin{note}
Оценка состоит из трёх слагаемых: оптимизационного --- порядка $\mathcal{O}(1/(\eta T))$, шумового --- порядка $\mathcal{O}(\eta\sigma^2)$ и дрейфового --- порядка $\mathcal{O}(\delta/\eta)$. При $\sigma = 0$ и $\delta \to 0$ восстанавливается скорость $\mathcal{O}(1/T)$ из теоремы \ref{theor:1}. Выбирая шаг
\begin{equation*}
\eta = \min\left\{ \frac{\alpha^2}{\Gamma\widetilde{L}},\; \alpha\sqrt{\frac{2\Delta_0}{\widetilde{L}\sigma^2 (T+1)}} \right\}
\end{equation*}
и $\beta$ так, что $1 - \beta = \mathcal{O}(\eta^2)$, получаем
\begin{equation*}
    \min_{0 \le t \le T} \mathbb{E}\|\nabla \widetilde{F}_t (w_t) \|^2 = \mathcal{O}\left( \frac{\Gamma^2 \widetilde{L} \Delta_0}{\alpha^2 (T+1)} + \frac{\Gamma}{\alpha}\sqrt{\frac{\Delta_0 \widetilde{L} \sigma^2}{T+1}} \right),
\end{equation*}
то есть скорость $\mathcal{O}(1/\sqrt{T})$, совпадающую по порядку со скоростью стохастического градиентного спуска в невыпуклом случае.
\end{note}

\begin{theorem}[стохастический сильно выпуклый случай]
\label{theor:4}
    Пусть выполняются условия теоремы \ref{theor:2}, в которых точный градиент заменён стохастическим, удовлетворяющим предположению \ref{ass:expectations}. Пусть по-прежнему
    \begin{equation*}
        0 < \eta \le \frac{\widetilde{\mu}\alpha^2}{\Gamma \widetilde{L}^2}, \qquad \rho_\beta \le \frac{\eta\widetilde{\mu}\alpha}{4\Gamma}.
    \end{equation*}
    Тогда для всех $T \ge 1$
    \begin{equation*}
        \mathbb{E} \left[ R_T^2 \right] \le \left(1 - \frac{\eta\widetilde{\mu}}{4\Gamma}\right)^T R_0^2 \;+\; \frac{8\Gamma\eta\sigma^2}{\widetilde{\mu}\alpha} \;+\; \frac{64\,\Gamma^3\lambda^2\Omega_0^2}{\eta^2\widetilde{\mu}^4}\,\rho_\beta^2.
    \end{equation*}
    В частности, для любого $\varepsilon > 0$, выбирая
    $\eta \le \min\left\{ \frac{\widetilde{\mu}\alpha^2}{\Gamma\widetilde{L}^2},\; \frac{\varepsilon\widetilde{\mu}\alpha}{24\Gamma\sigma^2} \right\}$ и $\beta$ так, что третье слагаемое не превосходит $\varepsilon/3$, получаем $\mathbb{E}[R_T^2] \le \varepsilon$ после
    $T = \frac{4\Gamma}{\eta\widetilde{\mu}}\log\frac{3R_0^2}{\varepsilon}$ итераций.
\end{theorem}

\begin{note}
Второе слагаемое --- классический <<шумовой шар>> стохастического градиентного метода с постоянным шагом: его радиус пропорционален $\eta\sigma^2$ и уменьшается вместе с шагом. Третье слагаемое --- дрейф цели $\widetilde{w}^*_t$, унаследованный из детерминированного случая (теорема \ref{theor:2}); он контролируется выбором $\beta \to 1$.
\end{note}

Эти теоремы устанавливают сходимость методов с предобуславливанием и затуханием весов при различных предположениях, а также определяют необходимое количество итераций для достижения заданной точности. Для наших задач сам факт сходимости этих методов имеет принципиальное значение.

Однако свойства решения $\widetilde{w}^*_t$ задачи
\begin{equation} \label{F_tilde_problem}
	\min_{w \in \mathbb{R}^d} \widetilde{F}_t(w) = f(w) + \widetilde{r}_t(w),
\end{equation}
к которому сходятся эти методы, требуют более глубокого исследования; ему посвящён следующий подраздел.
